% Generated draft: requires mathematical and source review.
\documentclass[11pt]{article}
\usepackage[margin=25mm]{geometry}
\usepackage{fontspec}
\setmainfont{Latin Modern Roman}
\usepackage{amsmath,amssymb,mathtools,unicode-math}
\setmathfont{Latin Modern Math}
\usepackage{enumitem,needspace,xurl,hyperref}
\setlength{\parindent}{0pt}
\setlength{\parskip}{5pt}
\setlength{\emergencystretch}{3em}
\newcommand{\N}{\mathbb{N}}
\newcommand{\Z}{\mathbb{Z}}
\newcommand{\Q}{\mathbb{Q}}
\newcommand{\R}{\mathbb{R}}
\newcommand{\C}{\mathbb{C}}
\newcommand{\F}{\mathbb{F}}
\newcommand{\A}{\mathbb{A}}
\newcommand{\PP}{\mathbb{P}}
\newcommand{\E}{\mathbb{E}}
\newcommand{\eps}{\varepsilon}
\newcommand{\dd}{\,\mathrm{d}}
\begin{document}
\section*{UnsolvedMath ID: 30004533}
\textbf{Source problem code:} catalogue entry
\subsection{Definitions and mathematical statement}
The supplied statement uses graded pieces but does not specify a ground field or its characteristic. To make the usual graded interpretation explicit, fix a field \(k\). Consider commutative standard graded algebras
\[
A=k[x_1,\ldots,x_r]/I=\bigoplus_{i\geq 0}A_i,
\]
where each variable has degree \(1\), \(I\) is a homogeneous ideal generated by quadratic forms, and \(A\) is finite-dimensional over \(k\). The codimension is \(\dim_k A_1=r\). Write \(\mathfrak m=\bigoplus_{i>0}A_i\) and define
\[
\operatorname{Soc}(A)=\{a\in A:\mathfrak m a=0\}.
\]
In this setting, the Artinian Gorenstein condition with socle degree \(4\) means
\[
\operatorname{Soc}(A)=A_4,\qquad \dim_k A_4=1,\qquad A_i=0\quad(i>4).
\]
Let \(QG_k(r)\) denote the class of these algebras. Whenever this class is nonempty, define
\[
\nu_k(r)=\min\{\dim_k A_2:A\in QG_k(r)\}.
\]
The supplied source writes these objects as \(QG(r)\) and \(\nu(r)\), without specifying \(k\). Its question is whether
\[
\lim_{r\to\infty}\frac{\nu(r)}{r^{2/3}}=6.
\]
Under the fixed-field interpretation above, this asks whether, for every \(\varepsilon>0\), there exists an integer \(R\geq 1\) such that for every integer \(r\geq R\),
\[
\left|\frac{\nu_k(r)}{r^{2/3}}-6\right|<\varepsilon.
\]
This formulation requires \(QG_k(r)\) to be nonempty for all sufficiently large integers \(r\). Neither that requirement nor the intended field hypotheses are supplied. The fixed-field notation clarifies a possible interpretation; it does not assert that the question is intended for every field or resolve the missing hypotheses.
\subsection{Short English statement}
Among quadratic-presented Artinian Gorenstein algebras of codimension \(r\) and socle degree \(4\), let \(\nu(r)\) be the smallest dimension of the degree-two component. The question asks whether \(\nu(r)/r^{2/3}\) tends to \(6\). The ground-field hypotheses and the range of codimensions for which the minimum is defined require clarification.
\subsection{Sources}
\begin{itemize}
\item Local UnsolvedMath registry, supplied record with permanent ID \(30004533\), titled \emph{Asymptotic Minimal Quadratic Growth of Gorenstein Algebras}; collection: Oberwolfach Reports Open Problems; exact source code: \texttt{catalogue entry}. No individual proposer is named. The supplied proposed year is \(2020\).
\item Supplied report citation: \emph{Lefschetz Properties in Algebra, Geometry and Combinatorics} (2021), Oberwolfach Reports, pp.\ 1539--1600, DOI 10.4171/owr/2020/31. \url{https://doi.org/10.4171/owr/2020/31}. The background describes Oberwolfach Reports as documenting mathematical workshops and their open-problem discussions.
\item Supplied literature reference, labeled ``Socle-degree-four study'': \url{https://www.sciencedirect.com/science/article/pii/S002240492030133X}.
\item Supplied literature reference, labeled ``Recent partial results'': \url{https://doi.org/10.1007/s13398-023-01523-6}.
\item Literature assessment dated exactly \texttt{2026-08-22}: the supplied curation labels the problem open and states that the asymptotic formula remains conjectural, with published partial results continuing to formulate the asymptotic question rather than proving the limit. This is supplied dated curation derived from the OWR record, not an independently verified literature assessment; its status and cited claims require independent verification.
\item The supplied statement status is \texttt{corrected\_verified}. Its verification note reports that the official report was checked and that the class \(QG(r)\), the minimum, and the conjectured limit were retained. This is a supplied verification claim, not a source consultation performed for this draft.
\end{itemize}
\end{document}
